\section{Coskeletal cutoffs and contractible trees}
\label{sec:cutoffs-trees}

We first establish two facts without using any elementary-expansion
assertion: higher cell attachments preserve lower cells, and free models
on globular sums are contractible.  We then fix a finite cutoff and set
up a descending induction.  Within one step of that induction, the
already established higher-dimensional expansions will suffice to
contract the tree diagrams needed for partial cylinders.

\subsection{The cutoff reflector}

Coskeletal globular models and their relation to truncated models are
studied in \cite[Section~2]{HenryLanari2023}.  We construct the reflector
needed here directly, keeping the chosen free order of the coherator.

For an integer $q\geq0$, call a model \emph{$q$-coskeletal} if every
globular boundary of dimension $j>q$ has exactly one filler.  Thus the
restriction map
\[
 \Hom(D_j,Y)\longrightarrow\Hom(\partial D_j,Y)
\]
is bijective for every $j>q$.  Let $\mathcal A_q$ be the full subcategory
of such models.  The uniqueness begins in dimension $q+1$; distinct
parallel $q$-cells are retained.
To recover a homotopy-class test in dimension $k$, we will choose
$q\geq k+1$, retaining both the tested cells and their joining homotopies.

\begin{lemma}[Cutoff reflector]\label{lem:cutoff}
The inclusion $\mathcal A_q\subseteq\mathcal A$ has a left adjoint
$R_q$.  Its unit $X\longrightarrow R_qX$ is bijective on the sets of
cells in every dimension at most $q$.  For $0\leq k<q$, the same
construction gives the reflector $\mathcal A_q\to\mathcal A_k$,
and $R_kR_q$ is naturally isomorphic to $R_k$.
\end{lemma}

\begin{proof}
On underlying globular sets, retain the cells and face maps of $X$
through dimension $q$.  In successive dimensions $j>q$, put one
$j$-cell over each parallel pair of $(j-1)$-cells.  Denote the resulting
globular set by $Z$.  There is a canonical globular map $\eta:X\to Z$,
equal to the identity through $q$.

To interpret the theory on $Z$, follow the \emph{given free cellular
presentation of $\mathcal C$ in its original order}.  Maintain that
$\eta$ respects the part of the theory already interpreted.  Suppose
a newly adjoined generating operation has output dimension $j$ and
arity $S$.  Its two prescribed faces are operations of the earlier
theory.

If $j\leq q$, the generating arity bound
$\dim(S)\leq j$ in Definition~\ref{def:coherator} makes the tuple sets
on $S$ in $X$ and $Z$ identical.  Use the original interpretation in
$X$.  Its prescribed face equations hold in $Z$: the tuple comes
from $X$, and $\eta$ already respects both earlier face operations.
It also respects the new operation, since it is the identity in
dimension $j$.

If $j>q$, the earlier interpretation gives a parallel boundary in
$Z$, and the new operation is its unique filler.  The image under
$\eta$ of the operation in $X$ has that boundary, so uniqueness again
shows compatibility with $\eta$.  At a limit stage take the union of
the preceding interpretations.  Freeness of the presentation now
gives a model structure on $Z$.  This argument places no bound on
the arities of \emph{composite} face operations and does not reorder
generators by dimension.

For a $q$-coskeletal model $Y$, a model map $f:X\to Y$ has a unique
globular factorization through $\eta$: use $f$ through dimension $q$
and unique fillers thereafter.  This factorization is a model map.
For a generator of output at most $q$, every tuple comes from $X$;
for a higher generator, preservation follows from uniqueness in
$Y$.  This proves the adjunction.  Starting with $X\in\mathcal A_q$
and performing the construction at $k<q$ proves the remaining
statements by the same universal property.
\end{proof}

Write $F_q=R_qF$ for the free $q$-coskeletal model functor.  Since
$R_q$ is a left adjoint, it preserves actual algebraic pushouts and
cellular presentations.  In $\mathcal A_q$ the images of $I_j$ and
$J_j$ are isomorphisms for $j>q$.  The image of $J_q$ freely adjoins
a parallel $q$-cell; its joining $(q+1)$-cell is forced by uniqueness.
This description includes all operations of the fixed theory on the
new cell.

Filtered colimits of $q$-coskeletal models, computed in
$\mathcal A$, remain $q$-coskeletal: both disks and their finite
boundaries commute with such colimits, so the required restriction
maps remain bijective.  Consequently filtered colimits in
$\mathcal A_q$ are computed on cells as well.  The adjunction for
$F_q$ now shows that $F_q(K)$ is finitely presentable whenever $K$
is a finite globular set.  The category $\mathcal A_q$ is
cocomplete by reflection, and the same ordinary small-object
construction is therefore available for its disk and boundary
maps.  No model structure is being used.

\begin{lemma}[Higher attachments preserve lower cells]
\label{lem:high-cells}
Fix $k\geq0$.  In $\mathcal A$, any relative cellular extension using
only $I_j$ with $j>k$ is bijective on all cells through dimension $k$.
The same holds in $\mathcal A_q$ whenever $k<q$.  Consequently these
conclusions also hold for relative extensions using only $J_j$
with $j>k$.
\end{lemma}

\begin{proof}
The reflection of $I_j$ into $\mathcal A_k$ is an isomorphism when
$j>k$: every $k$-coskeletal model has a unique filler of every
$j$-boundary.  Reflect a pushout, a coproduct of attachments, or a
transfinite composite of such maps.  Colimit preservation makes its
reflection an isomorphism.  Since both reflection units preserve
the cells through $k$, the original map is bijective on those cells.
Within $\mathcal A_q$, use the further reflection to $\mathcal A_k$.
Finally, $J_j$ is the composite of an $I_j$ attachment, adjoining its
new target $j$-cell, and an $I_{j+1}$ attachment, adjoining the
comparison cell.
\end{proof}

\subsection{Spines of globular sums}

We use an explicit recursive description of finite globular-sum
shapes.  For nonempty shapes $S_1,\ldots,S_p$, write
$[S_1,\ldots,S_p]$ for the shape with vertices
$v_0,\ldots,v_p$ whose cells over each consecutive pair
$(v_{i-1},v_i)$ form the suspension of $S_i$.  Suspension raises
dimensions by one and gives the suspended shape these two vertices
as poles.  Every shape is either a vertex or has this form.  To see
this from a table of dimensions, split its ordered amalgamation at
all gluings in dimension zero and lower every dimension in each
remaining block by one.  Repeating gives the stated recursion.

For a shape $S$, let $v(S)$ denote its first vertex.  Define a
subshape $T_m(S)$, for $m\geq0$, recursively.  A vertex is fixed, and
for $S=[S_1,\ldots,S_p]$ put
\begin{equation}\label{eq:spine-recursion}
\begin{split}
 T_0(S)&=[v(S_1),\ldots,v(S_p)],\\
 T_m(S)&=[T_{m-1}(S_1),\ldots,T_{m-1}(S_p)]\qquad(m\geq1).
\end{split}
\end{equation}
In the first line each $v(S_i)$ denotes the one-vertex subshape of
$S_i$.  Thus at depth $m$ the construction keeps every vertex and
one generating edge between consecutive vertices.  A branch ending
before that depth is retained in full.

\begin{lemma}[Finite spine]\label{lem:spine}
For every finite globular-sum shape $S$ and every $m\geq0$, the
subshape $T_m(S)$ has height at most $m+1$ and contains every cell of
$S$ through dimension $m$.  The inclusion $T_m(S)\to S$ is a finite
composite of pushouts of the globular source inclusions
$d_j\to d_{j+1}$ with $j>m$.
\end{lemma}

\begin{proof}
First, every shape $S$ can be built from $v(S)$ by finitely many
source inclusions $d_j\to d_{j+1}$ with $j\geq0$.  This is induction
on the recursive description.  If $S$ is a vertex there is nothing
to attach.  Otherwise, for $S=[S_1,\ldots,S_p]$, first build the
outer vertex string by source inclusions in dimension zero, using
in its $i$th slot the edge corresponding to $v(S_i)$.  Then build
each $S_i$ from $v(S_i)$ by the induction hypothesis and suspend its
attachments into that slot.  Under this suspension an elementary
pair of dimensions $(j,j+1)$ becomes one of dimensions $(j+1,j+2)$,
with exactly the same source-attachment description.  This follows
directly from the face maps: the two added cells acquire the fixed
outer poles, and their other faces are the suspended old faces.

The height and lower-cell claims for $T_m$ now follow directly from
Equation~\eqref{eq:spine-recursion}.  For the attachment claim when
$m=0$, rebuild each $S_i$ from $v(S_i)$ as above.  All these
attachments are suspended once, so their indices are at least
one.  For $m\geq1$, the induction hypothesis for
$T_{m-1}(S_i)\to S_i$ uses indices $j>m-1$; suspending gives indices
$j+1>m$.  Apply this construction in every slot.  There are finitely
many slots and attachments because $S$ is finite.
\end{proof}

Contractibility of free globular sums is established in
\cite[Proposition~3.8]{Lanari2020}.  The spine argument below proves
the statement directly in our convention and also gives its reflected form.

\begin{proposition}[Contractibility of free shapes]
\label{prop:shape-contractible}
For every nonempty finite globular-sum shape $S$, the models $F(S)$
and $F_q(S)$ are contractible, for every $q\geq0$.
\end{proposition}

\begin{proof}
Let $a,b$ be parallel $m$-cells of $F(S)$, where $m\geq0$ and in
dimension zero there is no parallelism condition.  By
Lemmas~\ref{lem:spine} and~\ref{lem:high-cells}, the map
$F(T_m(S))\to F(S)$ is bijective on every cell set through $m$.
Hence $a,b$ have unique lifts $a',b'$ there, and the lifts are
parallel: their lower faces agree by the same injectivity.

The free model on a globular-sum shape is represented by that
object of $\mathcal C$.  Thus $a',b'$ are a parallel pair of theory
arrows with domain the $m$-globe and codomain $T_m(S)$.
Since $\dim(T_m(S))\leq m+1$, this pair is admissible for a filler
of dimension $m+1$.  Contractibility of $\mathcal C$ supplies such a
filler.  Its image joins $a$ to $b$ in $F(S)$.
The model is nonempty because $S$ has a vertex, so the
contractibility criterion of Proposition~\ref{prop:weak-test}
applies.  No assertion about weak equivalences of elementary
attachments was used here.

For $F_q(S)$, a parallel pair of $m$-cells with $m<q$ lifts to
$F(S)$ by Lemma~\ref{lem:cutoff}; the filler just constructed is
preserved through dimension $m+1\leq q$.  If $m\geq q$, the
required filler has dimension $m+1>q$ and exists uniquely.  The
vertex is retained, so $F_q(S)$ is also contractible.
\end{proof}

\subsection{The descending induction and its upper-cell consequence}

Fix $q\geq0$.  Until the cutoff is removed in the next section,
all constructions take place in $\mathcal A_q$.  We keep the notation
$F_q$ for free models, and write $D_j,\partial D_j,I_j,J_j$ for
their images under $R_q$.  Weak equivalence means the same
component and based-homotopy-group test on these models; all the
properties in Section~\ref{sec:preliminaries} remain applicable.

For $0\leq n\leq q$, let $\mathsf E(q,n)$ be the following assertion:
\begin{quote}
For every cellular object $X\in\mathcal A_q$ and every model map
$a:D_n\to X$, the actual pushout of $J_n$ along $a$ is a weak
equivalence.
\end{quote}
Here cellular objects and extensions may have any set-sized
transfinite presentation.  The attaching map $a$ need not select
a presentation generator.  We will prove $\mathsf E(q,n)$ by
descending induction on $n$.

\begin{lemma}[Top-dimensional base case]\label{lem:cutoff-base}
The assertion $\mathsf E(q,q)$ holds for every $q\geq0$.
\end{lemma}

\begin{proof}
Let $i:X\to Y$ be a pushout of $J_q$.  It has a retraction
$r:Y\to X$, obtained by sending the new $q$-cell to the attaching
cell and its comparison to a unit.  For $q>0$,
Lemma~\ref{lem:high-cells} shows that $i$ is bijective on cells
through dimension $q-1$.

Check the exact-boundary criterion of
Proposition~\ref{prop:weak-test} in dimension $d$.  If $d<q$, the
given $d$-cell of $Y$ comes from a unique $d$-cell of $X$, whose
boundary is the prescribed one by injectivity on lower cells.
Its image is the given cell, so its unit supplies the homotopy.
If $d=q$, apply $r$ to the given cell.  Its boundary is the
prescribed boundary because $ri=\id_X$, and its image is parallel
to the given cell.  These two $q$-cells have a unique joining
$(q+1)$-cell.  If $d>q$, the prescribed source boundary has a
unique filler in $X$, whose image is the given filler in $Y$.
This proves the criterion in every dimension.
For $q=0$, the first case is absent; the same retraction and
unique-arrow argument proves the dimension-zero case directly.
\end{proof}

For the rest of this section and the construction in the next,
fix $0\leq n<q$ and assume $\mathsf E(q,j)$ for every
$n<j\leq q$.  Call an index \emph{upper} if it is strictly greater
than $n$.  Thus an upper $I$-cellular extension uses only $I_j$
with $j>n$, and likewise for upper $J$-cellular extensions.
Every upper $J$-cellular extension with cellular source is in
$\mathcal W$: the successor stages use the induction hypothesis
(or an isomorphism for $j>q$), and limit stages use
Lemma~\ref{lem:weak-closure}.  Every intermediate object remains
cellular, since each $J_j$ is a composite of two boundary
attachments.

\begin{lemma}[Upper-cell splitting and pushout]
\label{lem:upper-cell}
Suppose $i:U\to V$ is an upper $I$-cellular extension, $U$ is
cellular, and $i\in\mathcal W$.  Then $i$ admits a retraction.
For every map $U\to U'$ with $U'$ cellular, the pushout
\[
 U'\longrightarrow U'\amalg_U V
\]
also belongs to $\mathcal W$.
\end{lemma}

\begin{proof}
Apply the ordinary small-object construction with the upper
$J_j$ to factor $i$ as
\[
 U\xrightarrow{\,a\,}Z\xrightarrow{\,p\,}V,
\]
where $a$ is upper $J$-cellular and $p$ lifts against every upper
$J_j$.  The preceding induction hypotheses imply $a\in\mathcal W$,
so $p\in\mathcal W$ by two-out-of-three.  By
Lemma~\ref{lem:lifting-test}, $p$ lifts against every upper $I_j$:
the approximate lift of a $j$-cell is corrected by lifting its
$(j+1)$-dimensional homotopy against $J_j$.  In particular the
lowest index required is $J_{n+1}$, which is already available.

Since $i$ is upper $I$-cellular, there is a map $s:V\to Z$ with
$si=a$ and $ps=\id_V$.  Thus, in the category under $U$, the map
$i$ is a retract of $a$.
Each individual $J_j$ attachment admits a unit retraction.  A
transfinite composite of such attachments also retracts to its
initial object: choose the maps back recursively at successor
stages, and use their compatible colimit map at a limit stage.
Let $t:Z\to U$ be the resulting retraction of $a$.  Then $ts$ is
a retraction of $i$.

Push the retract diagram out along $U\to U'$.  The pushout of
$i$ is a retract of the pushout of $a$, and the latter is upper
$J$-cellular with cellular source $U'$.  It is therefore in
$\mathcal W$ by the induction hypotheses.  Retract closure proves
the claim.
\end{proof}

\subsection{Trees in a fixed globular dimension}

Define a \emph{tree diagram at depth $n$} as follows.  At depth zero,
it is a finite nonempty oriented tree, regarded as a globular set
of dimensions zero and one; a single vertex is allowed.  At
depth $n>0$, it is either a vertex or a string
$[K_1,\ldots,K_p]$ with every $K_i$ a tree diagram at depth $n-1$.
Such a diagram $K$ is a finite globular set of height at most
$n+1$.  In each depth-$n$ slot, the tree vertices are $n$-cells
with one common lower boundary, and its edges are $(n+1)$-cells.
All orientations of the edges are allowed.  At $n=0$ the lower
boundary condition is empty.

The free model $F_q(K)$ is cellular: attach its finitely many
globular cells along their boundaries in increasing dimension and
apply the free functor.  Its finite generating globular set must
be distinguished from the generally infinite collection of cells
created by the operations of $\mathcal C$.

\begin{fixedfigure}
\centering
\begin{tikzpicture}[>=Stealth,
    vertex/.style={circle,fill=black,inner sep=1.55pt},
    every node/.style={font=\small}]
\begin{scope}
\node[vertex,label=above:\(x_0\)] (x0) at (0,1.2) {};
\node[vertex,label=above:\(x_1\)] (x1) at (1.45,1.2) {};
\node[vertex,label=above:\(x_2\)] (x2) at (2.9,1.2) {};
\node[vertex,label=below:\(y_0\)] (y0) at (0,0) {};
\node[vertex,label=below:\(y_1\)] (y1) at (1.45,0) {};
\node[vertex,label=below:\(y_2\)] (y2) at (2.9,0) {};
\draw[->] (x0) -- (x1);
\draw[->] (x1) -- (x2);
\draw[->] (x0) -- (y0);
\draw[->] (x1) -- (y1);
\draw[->] (x2) -- (y2);
\node at (1.45,-1.05) {a tree slot};
\end{scope}
\begin{scope}[xshift=4.4cm]
\node[vertex,label=above:\(y_0\)] (z0) at (0,0.6) {};
\node[vertex,label=above:\(x_0\)] (z1) at (1,0.6) {};
\node[vertex,label=above:\(x_1\)] (z2) at (2,0.6) {};
\node[vertex,label=above:\(y_1\)] (z3) at (3,0.6) {};
\node[vertex,label=above:\(x_2\)] (z4) at (4,0.6) {};
\node[vertex,label=above:\(y_2\)] (z5) at (5,0.6) {};
\draw[->] (z0) -- (z1);
\draw[->] (z1) -- (z2);
\draw[->] (z2) -- (z3);
\draw[->] (z3) -- (z4);
\draw[->] (z4) -- (z5);
\node at (2.5,-1.05) {a string on the same vertices};
\end{scope}
\end{tikzpicture}
\caption{A tree in one depth-\(n\) slot and a string on the same vertices.
Vertices are \(n\)-cells and arrows are
\((n+1)\)-cells.  The comparison maps between the free models fix the
named vertices and send generating edges to composites along connecting
paths, using inverses where necessary.  The \(J_{n+1}\)-duplication
argument in Proposition~\ref{prop:tree-contractible} promotes the
resulting edge homotopies to a statement about the full models.}
\label{fig:tree-string}
\end{fixedfigure}


\begin{proposition}[Contractibility of tree diagrams]
\label{prop:tree-contractible}
Under the fixed-cutoff induction hypotheses above, $F_q(K)$ is
contractible for every tree diagram $K$ at depth $n$.
\end{proposition}

\begin{proof}
In each depth-$n$ slot, order its tree vertices in any way and
replace the tree by the forward linear string on those same
labeled vertices.  A singleton remains a singleton.  The resulting
globular set $K^{\mathrm{lin}}$ is a globular-sum shape, and its
cells through dimension $n$ are canonically the same as those of
$K$; Figure~\ref{fig:tree-string} illustrates one slot.  Reordering
within a slot leaves unchanged any outer poles and
every earlier concatenation or suspension.  Put
\[
 A=F_q(K),\qquad L=F_q(K^{\mathrm{lin}}).
\]
The model $L$ is contractible by
Proposition~\ref{prop:shape-contractible}.

There are model maps
\[
 f:A\longrightarrow L,\qquad g:L\longrightarrow A
\]
fixing the generating cells through dimension $n$.  To define
$f$, send each tree edge to a composite along the path between
its endpoints in the corresponding line, using inversion for
any backward traversal.  Define $g$ using paths in the tree.
All compositions are at level $n$.  Within one slot the vertices
are parallel $n$-cells, so these assignments have exactly the
required faces.  They extend to model maps by the free universal
property.

For each $(n+1)$-dimensional generating edge $e$ of $K$, choose
a cell
\begin{equation}\label{eq:tree-generator-homotopy}
 \alpha_e:e\longrightarrow gf(e)
 \qquad\text{of dimension }n+2\text{ in }A.
\end{equation}
Such a cell exists by the groupoid calculus of
Section~\ref{sec:preliminaries}.  Indeed, $gf(e)$ represents a
walk in the original tree with the same endpoints as $e$.
Here an inverse of a composite is expanded at the level of homotopy
classes, using uniqueness of inverses in the hom-set groupoid.
Successively delete backtracking pairs from this walk.  A walk
in a tree without backtracking is the unique simple path between
its endpoints, here the edge $e$.  Associativity, units, and
inverse cancellation hold for $(n+1)$-cells modulo
$(n+2)$-cell homotopy, so the two cells represent the same class.
The homotopy relation is given by existence of a single joining
cell, giving Equation~\eqref{eq:tree-generator-homotopy}.

We next turn these homotopies on generators into a statement
about the whole model.  For every such edge $e$, adjoin a
duplicate $e'$ with the same boundary and a cell
$\delta_e:e\to e'$ of dimension $n+2$.  Let $Q$ be the resulting
model and $a_0:A\to Q$ its canonical map.  These are finitely
many $J_{n+1}$ attachments, so $a_0\in\mathcal W$ by the already
proved higher-index case.  Define a second model map
$a_1:A\to Q$ by fixing every generator through dimension $n$
and sending $e$ to $e'$.  This is permitted because
$A=F_q(K)$ and $K$ has no generators above dimension $n+1$.
The face equations of this globular assignment are the only
conditions needed to extend it to the full model.

There is a common retraction $r:Q\to A$: keep the original
generators, send $e'$ to $e$, and send $\delta_e$ to a unit on
$e$.  Thus $ra_0=ra_1=\id_A$.  Two-out-of-three gives first
$r\in\mathcal W$ and then $a_1\in\mathcal W$.
Likewise Equation~\eqref{eq:tree-generator-homotopy} defines a
model map $h:Q\to A$ which is the identity on the original copy,
sends $e'$ to $gf(e)$, and sends $\delta_e$ to $\alpha_e$.
The free universal property gives
\[
 ha_0=\id_A,\qquad ha_1=gf.
\]
It follows that $h\in\mathcal W$ and $gf\in\mathcal W$.
This promotion uses only $J_{n+1}$, not $J_n$.

The map $gf$ factors through the contractible model $L$.  Hence
its map on components factors through a singleton and its maps
on positive based homotopy groups factor through trivial groups.
Since these maps are respectively bijective and isomorphisms,
$A$ has one component and trivial homotopy groups.  It is
nonempty, so it is contractible.

When $q=n+1$, the attachments used above are precisely $J_q$,
whose case was proved in Lemma~\ref{lem:cutoff-base}.  They retain
the duplicate $q$-generators, while the joining $(q+1)$-cells
are unique.  The definitions of $a_1,r,h$ still have exactly
their stated universal properties.  Thus this boundary case
requires no additional induction step.
\end{proof}

\subsection{Whiskers on globular sums}

Let $S$ be a finite globular-sum shape with $\dim(S)\leq q$.
For $\epsilon\in\{0,1\}$, define $H_\epsilon(S)$ by starting
with $F_q(S)$ and, for each actual $n$-cell $x$ of the globular
set $S$, adjoining a duplicate $x'$ with the same boundary and
one $(n+1)$-cell joining them.  Orient the join $x\to x'$ when
$\epsilon=0$ and $x'\to x$ when $\epsilon=1$.  Every original
higher generating cell of $S$ keeps its original boundary.
These are cellular constructions: attach each duplicate along
its boundary, then its joining cell.  For $n=0$ the duplicate
attachment is an attachment of a vertex.  If $S$ has no
$n$-cells, set $H_\epsilon(S)=F_q(S)$.

\begin{proposition}[Contractibility after adjoining whiskers]
\label{prop:whiskered-shape}
Under the fixed-cutoff induction hypotheses above,
$H_\epsilon(S)$ is contractible for both choices of $\epsilon$.
The same conclusion holds if the orientation of the join is
chosen independently for each $n$-cell of $S$.
\end{proposition}

\begin{proof}
Set $S'=T_n(S)$.  It contains exactly the same $n$-cells as $S$,
so the same duplicates and joins can be adjoined to either
shape.  In $S'$ every slot at depth $n$ is a linear string or a
singleton.  Adding one leaf and its edge at each vertex turns
each such slot into a finite oriented tree.  Thus
$H_\epsilon(S')$ is the free model on a tree diagram at depth
$n$, and is contractible by
Proposition~\ref{prop:tree-contractible}.  The same argument
allows independent orientations.

By Lemma~\ref{lem:spine}, $S'\to S$ is a finite composite of
source attachments with indices strictly greater than $n$.
Those attachments add no further $n$-cells.  Perform them after
adjoining the fixed set of duplicates and joins.  Their original
attaching maps use the retained generators of $S$ and their
original faces, so the resulting model is exactly
$H_\epsilon(S)$ by the pushout universal properties.  Therefore
$H_\epsilon(S')\to H_\epsilon(S)$ is upper $J$-cellular.  Its
source and all intermediate objects are cellular, so the
descending induction hypotheses make it a weak equivalence.
Contractibility follows from that of $H_\epsilon(S')$.
If the recursion ends before depth $n$, there are no joins in
that branch and it remains unchanged throughout.
\end{proof}

The hypotheses $0\leq n<q$ and $\mathsf E(q,j)$ for $j>n$
remain in force in the next section.  The tree and whisker
propositions have established the needed non-pasting diagrams
using only those hypotheses; the assertion $\mathsf E(q,n)$
itself is still to be proved.
